\begin{figure}[H]
    \centering
    \begin{tikzpicture}[x=1.35cm, y=0.85cm]

        \draw[flow] (0,0) -- (7.6,0);
        \draw[flow] (0,0) -- (0,6.4);
        \node[note, font=\scriptsize, anchor=north east] at (7.6,-0.12)
            {υπολογιστική ένταση $I = W/Q$};
        \node[note, font=\scriptsize, anchor=south west, rotate=90] at (-0.42,0)
            {επιτεύξιμη επίδοση $P$};

        \draw[line width=0.8pt, draw=figaieD] (0.4,1.4) -- (4.0,5.0) -- (7.2,5.0);
        \draw[line width=0.5pt, draw=figcacD] (0.4,2.4) -- (3.0,5.0);
        \draw[line width=0.5pt, draw=figregD] (0.4,3.4) -- (2.0,5.0);

        \draw[dashed, draw=figmute, line width=0.4pt] (0,5.0) -- (7.2,5.0);
        \node[notemute, font=\tiny, anchor=east] at (-0.06,5.0) {$\pi$};

        \node[notemute, font=\tiny, anchor=south west, text=figregD] at (0.34,3.44) {\tl{L2}};
        \node[notemute, font=\tiny, anchor=south west, text=figcacD] at (0.34,2.44) {\tl{L3}};
        \node[notemute, font=\tiny, anchor=south west, text=figaieD] at (0.34,1.44) {\tl{DRAM}};

        \fill[figink] (2.0,3.0) circle (2.2pt);
        \node[note, font=\tiny, anchor=north east] at (1.94,2.94) {\tl{A}};
        \fill[figink] (5.6,5.0) circle (2.2pt);
        \node[note, font=\tiny, anchor=south] at (5.6,5.14) {\tl{B}};
        \fill[figink] (3.6,2.0) circle (2.2pt);
        \node[note, font=\tiny, anchor=north east] at (3.54,1.94) {\tl{C}};

        \node[notemute, font=\tiny, anchor=west, align=left] at (2.20,2.85)
            {\tl{A}: φραγμένο από τη μνήμη};
        \node[notemute, font=\tiny, anchor=east, align=right] at (5.42,4.62)
            {\tl{B}: φραγμένο από την\\ υπολογιστική ικανότητα};
        \node[notemute, font=\tiny, anchor=west, align=left] at (3.80,1.85)
            {\tl{C}: φραγμένο από καθυστερήσεις\\ ή ανεπαρκή παραλληλία};

        \node[note, font=\tiny, anchor=north west, align=left] at (0.0,-0.62)
            {Λογαριθμική κλίμακα και στους δύο άξονες.};

    \end{tikzpicture}
    \caption{Η ανάγνωση του μοντέλου \tl{roofline}. Το κεκλιμένο τμήμα αντιστοιχεί
    στο όριο $\beta I$ που επιβάλλει το εύρος ζώνης και το οριζόντιο στο όριο $\pi$
    που επιβάλλει η υπολογιστική ικανότητα· η επιτεύξιμη επίδοση είναι το ελάχιστο
    των δύο. Επειδή το εύρος ζώνης εξαρτάται από το μέγεθος του συνόλου εργασίας,
    το κεκλιμένο τμήμα δεν είναι μία ευθεία αλλά μια κλίμακα από παράλληλες
    ευθείες, μία ανά βαθμίδα της ιεραρχίας μνήμης, και κάθε μέτρηση συγκρίνεται με
    εκείνη που αντιστοιχεί στο δικό της σύνολο εργασίας. Το σημείο \tl{A}
    περιορίζεται από τη μνήμη και θα ωφεληθεί μόνο από μείωση της κίνησης
    δεδομένων, το \tl{B} από την υπολογιστική ικανότητα, ενώ το \tl{C} απέχει και
    από τα δύο ταβάνια: εκεί ο περιορισμός δεν είναι το υλικό αλλά η οργάνωση της
    εκτέλεσης.}
    \label{fig:meth-roofline}
\end{figure}
